%! Parametrically Defined Circles and Lines — Unit 8, Lesson 8.3 — Guided Notes (Answer Key)
\documentclass[10pt]{article}
\usepackage{precalculus-article}
\usepackage{precalculus-key}   % pulls in -boxes; do NOT also load -boxes
\newcommand{\TallMath}[1]{$\displaystyle #1\rule[-1.4em]{0pt}{3.2em}$}

\begin{document}

\pageheader{Unit 8, Lesson 8.3}{Guided Notes --- Answer Key}
\namedateperiod

\vspace{0.08in}

%----------------------------------------------------------------------
% Objectives
%----------------------------------------------------------------------
\begin{objectivebox}
\textbf{I will be able to\ldots}
\begin{enumerate}[leftmargin=1.5em, itemsep=4pt]
  \item \textbf{LO 4.4.A} --- Express the motion of a point around a
    circle parametrically by applying transformations to
    $(x(t),\,y(t)) = (\cos t,\,\sin t)$, and parametrize the segment
    between two given points using an initial position and constant
    rates of change. \ansline{see notes boxes below}
\end{enumerate}
\end{objectivebox}

\vspace{0.08in}

%----------------------------------------------------------------------
% Vocabulary
%----------------------------------------------------------------------
\begin{vocabbox}
\begin{description}[leftmargin=0pt, itemsep=6pt]
  \item[\termblanklong{unit circle parametrization}]
    The equations $x(t) = \cos t$, $y(t) = \sin t$, $0 \le t \le 2\pi$
    model one full \ans{counterclockwise} revolution starting at
    \ans{$(1,0)$} and centered at the \ans{origin}.

  \item[\termblanklong{radius / center}]
    In $(h + r\cos t,\; k + r\sin t)$: $r$ is the
    \ans{radius} and $(h,\,k)$ is the
    \ans{center} of the circular path.

  \item[\termblanklong{orientation}]
    Using $(\cos t,\,\sin t)$ gives \ans{counterclockwise} motion.
    Replacing $t$ with $-t$ gives \ans{clockwise} motion.

  \item[\termblanklong{linear parametrization}]
    The parametrization $x(t) = x_1 + (x_2-x_1)t$,
    $y(t) = y_1 + (y_2-y_1)t$, $0\le t\le 1$, traces the segment from
    \ans{$(x_1,y_1)$} to \ans{$(x_2,y_2)$}.
\end{description}
\end{vocabbox}

\vspace{0.08in}

%----------------------------------------------------------------------
% Hook
%----------------------------------------------------------------------
\begin{hookbox}
\textbf{Entry Question:} A Ferris wheel has radius 30 ft with its center
35 ft above the ground. A rider boards at the rightmost seat.

\vspace{4pt}
\textbf{Q1.} What are the rider's starting coordinates $(x,\,y)$ relative
to a center at the origin?

\vspace{4pt}
\ans{$(30,\;0)$} relative to the center; in ground coordinates, $(30,\;35)$.

\textbf{Q2.} As the wheel turns counterclockwise, what do you expect
to happen to the rider's height $y$?

\vspace{4pt}
\ans{The height increases to a maximum of $35 + 30 = 65$ ft (at the top) then decreases back to $35 - 30 = 5$ ft (at the bottom), following a sine curve.}
\end{hookbox}

\vspace{0.08in}

%----------------------------------------------------------------------
% Notes Box 1 — Unit Circle
%----------------------------------------------------------------------
\begin{notesbox}{LO 4.4.A --- Part 1: Unit Circle Parametrization (EK 4.4.A.1)}

\textbf{Definition:} The parametric function
$(x(t),\,y(t)) = (\cos t,\,\sin t)$, $0 \le t \le 2\pi$, traces the
\ans{unit} circle exactly \ans{once}, moving \ans{counterclockwise}.

\vspace{6pt}
\textbf{Complete the table:}

\vspace{4pt}
\begin{center}
\renewcommand{\arraystretch}{1.7}
\begin{tabular}{|c|c|c|c|}
\hline
$t$ & $x(t) = \cos t$ & $y(t) = \sin t$ & $(x,\;y)$ \\
\hline
$0$      & \ans{$1$}  & \ans{$0$}  & \ans{$(1,\;0)$}  \\
\hline
$\pi/2$  & \ans{$0$}  & \ans{$1$}  & \ans{$(0,\;1)$}  \\
\hline
$\pi$    & \ans{$-1$} & \ans{$0$}  & \ans{$(-1,\;0)$} \\
\hline
$3\pi/2$ & \ans{$0$}  & \ans{$-1$} & \ans{$(0,\;-1)$} \\
\hline
$2\pi$   & \ans{$1$}  & \ans{$0$}  & \ans{$(1,\;0)$}  \\
\hline
\end{tabular}
\end{center}

\begin{teachernote}
  Emphasize that the table values come directly from the unit-circle
  cosine/sine values from the warm-up. Students can verify geometrically
  by looking at the standard unit circle diagram.
\end{teachernote}

\end{notesbox}

\vspace{0.08in}

%----------------------------------------------------------------------
% Notes Box 2 — Transformed Circles
%----------------------------------------------------------------------
\begin{notesbox}{LO 4.4.A --- Part 2: Transformations of the Circle (EK 4.4.A.2)}

\textbf{General form:}
\[
  x(t) = \ans{h} + \ans{r}\cos t, \qquad
  y(t) = \ans{k} + \ans{r}\sin t, \qquad
  0 \le t \le 2\pi
\]
\begin{itemize}[leftmargin=1.4em, itemsep=2pt]
  \item Center: \ans{$(h,\,k)$} \quad Radius: \ans{$r$}
  \item Counterclockwise when $t$ increases \quad / \quad Clockwise: replace $t$
    with \ans{$-t$}
\end{itemize}

\vspace{8pt}
\textbf{Ferris Wheel Example.} Radius 30 ft, center at $(0,\,35)$, counterclockwise:

\vspace{4pt}
$x(t) = $ \ans{$30\cos t$} \qquad $y(t) = $ \ans{$35 + 30\sin t$}

\vspace{6pt}
Evaluate at $t = \pi/2$: \quad $x(\pi/2) = $ \ans{$0$}, \quad $y(\pi/2) = $ \ans{$65$}

The rider is at position \ans{$(0,\;65)$} after a quarter turn.

\vspace{6pt}
\textbf{Clockwise version} (replace $t$ with $-t$):

$x(t) = $ \ans{$30\cos(-t) = 30\cos t$} \qquad $y(t) = $ \ans{$35 + 30\sin(-t) = 35 - 30\sin t$}

\begin{teachernote}
  Common error: students negate \emph{both} components for clockwise.
  Stress that $\cos(-t) = \cos t$ (cosine is even), so only $\sin$ changes sign.
\end{teachernote}

\end{notesbox}

\vspace{0.08in}

%----------------------------------------------------------------------
% Notes Box 3 — Linear Segment
%----------------------------------------------------------------------
\begin{notesbox}{LO 4.4.A --- Part 3: Linear Parametrization (EK 4.4.A.3)}

\textbf{Segment from $(x_1,\,y_1)$ to $(x_2,\,y_2)$:}
\[
  x(t) = x_1 + (x_2-x_1)\,t, \qquad y(t) = y_1 + (y_2-y_1)\,t, \qquad 0 \le t \le 1
\]
At $t=0$: position is \ans{$(x_1,\,y_1)$}.
At $t=1$: position is \ans{$(x_2,\,y_2)$}.
At $t=0.5$: position is the \ans{midpoint}.

\vspace{6pt}
\textbf{Example:} Segment from $(2,\,5)$ to $(6,\,1)$.

\vspace{4pt}
$x(t) = 2 + \bigl($ \ans{$4$} $\bigr)t = $ \ans{$2 + 4t$} \qquad
$y(t) = 5 + \bigl($ \ans{$-4$} $\bigr)t = $ \ans{$5 - 4t$}

\vspace{6pt}
\textbf{Complete the table:}

\vspace{4pt}
\begin{center}
\renewcommand{\arraystretch}{1.6}
\begin{tabular}{|c|c|c|c|}
\hline
$t$ & $x(t)$ & $y(t)$ & $(x,\;y)$ \\
\hline
$0$   & \ans{$2$} & \ans{$5$} & \ans{$(2,\;5)$} \\
\hline
$0.5$ & \ans{$4$} & \ans{$3$} & \ans{$(4,\;3)$} \\
\hline
$1$   & \ans{$6$} & \ans{$1$} & \ans{$(6,\;1)$} \\
\hline
\end{tabular}
\end{center}

\end{notesbox}

\end{document}
